This paper uses a project from MIT's System Design and Management program to examine how clinical AI should be designed as a managed system, not as a standalone model. The case, Sanguine, is a physician-facing interpretation layer for Mayo Clinic laboratory results. Across the project, the team moved from stakeholder analysis and problem framing to architecture decisions, tradespace exploration, uncertainty modeling, development planning, and risk mitigation. The main finding is practical: in clinical AI, model capability is only one part of the design problem. Evidence grounding, release control, workflow fit, regulatory readiness, and institutional trust can dominate the final architecture and the project baseline. The resulting recommendation favors a higher-cost, risk-mitigated plan because the lower-cost alternatives preserve risks that could block clinical validation or damage trust.
